\chapter{Conclusion}

\section{Summary of Key Findings}
This study assessed the resilience of the transport network in Southern Italy under seismic disruption scenarios, focusing on multiple graph theory indicators to evaluate network performance before and after disruptions. The research identified critical links and nodes, connectivity shifts, and redundancy changes, providing a comprehensive understanding of how different seismic faults impact regional mobility.

\begin{itemize}
    \item \textbf{Betweenness Centrality Analysis} revealed that disruptions shift network importance, concentrating traffic on alternative routes, particularly in Campania and Puglia. The most vulnerable links were motorways and primary roads, emphasizing the importance of maintaining these corridors for network resilience.
    
    \item \textbf{Node Centrality Analysis} showed that intersections near Seismic Fault 1 and Fault 2 played a crucial role in overall connectivity, with high-intensity earthquakes significantly increasing nodal vulnerability.
    
    \item \textbf{Beta Index and Gamma Index} computations highlighted that network connectivity declined significantly after disruptions, especially in high-intensity earthquake scenarios, where many road segments became isolated.
    
    \item \textbf{Clustering Coefficient and Average Degree (CD) Analysis} indicated that local connectivity fragmentation was more pronounced in rural areas, whereas urban areas showed better resilience due to alternative routes.
    
    \item \textbf{Redundancy Analysis} revealed that high-intensity earthquakes caused the complete isolation of several zones, particularly in eastern Campania, where limited alternative routes exacerbated network fragmentation. The A16 motorway was identified as a critical transport corridor, and its disruption led to significant travel time increases and severe accessibility loss, further highlighting the region’s vulnerability to seismic events.
\end{itemize}

The results underscore the necessity of strategic interventions to enhance network robustness against seismic events:

\begin{itemize}
    \item Strengthening critical road corridors, especially motorways and trunk roads, to prevent network fragmentation.
    \item Enhancing redundancy measures in Campania and Calabria is essential, as these regions have fewer alternative routes, making them particularly susceptible to severe disruptions. Strengthening key road links and improving network adaptability can help mitigate the impact of seismic-induced failures and maintain regional mobility.
    \item Enhancing emergency response planning by identifying and prioritizing high-risk links and nodes to minimize congestion and facilitate recovery.
    \item Developing alternative expressways and bypass routes in eastern Campania is crucial to ensuring that seismic disruptions do not severely affect the broader transport network. Strategic investments in new transport corridors and reinforcement of existing infrastructure can enhance regional connectivity and reduce dependency on single critical routes, thereby improving overall network resilience.Utilizing GIS-based and graph-theoretic models for ongoing seismic risk assessment and adaptive infrastructure planning.
    
\end{itemize}

\section{Contributions of the Study}
This research integrates network science and resilience analysis to offer a quantitative evaluation of seismic impacts on transport systems. The study’s main contributions include:

\begin{itemize}
    \item Applying graph theory metrics to evaluate transport network resilience in an interregional context.
    \item Incorporating GIS-based spatial analysis to map and visualize post-disruption connectivity changes.
    \item Providing a structured framework for assessing and improving transport resilience, adaptable to other seismically active regions.
\end{itemize}

\section{Limitations and Future Research}
While this study provides valuable insights, some limitations remain:

\begin{itemize}
    \item The model assumes uniform road disruption probabilities, whereas real-world failures may depend on soil conditions, structural integrity, and traffic loads.
    \item The analysis focuses on road transport, but future studies could incorporate multimodal networks (rail, maritime, air) to assess comprehensive mobility resilience.
    \item Long-term recovery dynamics were not included; future research should explore how disrupted networks gradually regain functionality over time.
    \item Climate-induced disruptions (e.g., floods, landslides) should be considered alongside seismic risks to create a multi-hazard resilience framework.
\end{itemize}

\section{Final Remarks}
This research highlights the critical need for resilient transport infrastructure in Southern Italy, where seismic events pose a significant threat to regional connectivity and economic stability. By combining network analysis, GIS methodologies, and resilience indicators, this study provides actionable recommendations to enhance transport robustness, mitigate disruptions, and support long-term planning efforts. The findings contribute to the broader discourse on transport resilience, emphasizing the role of data-driven decision-making in safeguarding mobility in disaster-prone regions.
